\documentclass[11pt]{article}
\usepackage[margin=1in]{geometry}
\usepackage{amsmath,amssymb,amsthm,enumitem,hyperref}
\hypersetup{colorlinks=true,linkcolor=blue,urlcolor=blue,citecolor=blue}
\newtheorem{theorem}{Theorem}
\newcommand{\LL}{\mathcal L}
\title{Hoggatt sums fail infinite log-convexity for every $m\ge4$\\\large MathDB 368345}
\author{Principia Math In Progress}
\date{September 18, 2026}
\begin{document}
\maketitle
\begin{abstract}
We disprove Conjecture 4.3 of Fang--Zhang--Zhao in its current
arXiv formulation. For parameter $m=4$ or $5$, the fourth
log-convexity iterate has a negative entry. For every $m\ge6$,
the third iterate already has a negative entry. The separate Baxter
case $m=3$ is not decided. The identified journal final has not
been read; campaign eligibility remains pending.
\end{abstract}

\section*{Task 1. Full disproof and uniform extension}
Use exactly the source's convention
\[
 (\LL a)_n=a_na_{n+2}-a_{n+1}^2.
\]
Infinite log-convexity requires every entry of $\LL^r a$ to be
nonnegative for every integer $r\ge1$. Define the Hoggatt sums
from \cite[p.~3]{FZZ} by
\[
 A_n(m)=H_n^{[m]}(1,1)=
 \sum_{k=0}^{n-1}
 \frac{\prod_{i=0}^{m-1}\binom{n+m-2}{k+i}}
 {\prod_{j=1}^{m-1}\binom{n+m-2}{j}},\qquad n\ge1.       \tag{1}
\]
Conjecture 4.3 asserts infinite log-convexity for every fixed
integer $m\ge2$.

\begin{theorem}
The sequence $(A_n(m))_{n\ge1}$ is not infinitely log-convex
for any integer $m\ge4$. More precisely,
\[
 \begin{aligned}
 (\LL^4 A(4))_2&=-111509459048444462640611465400686603136,\\
 (\LL^4 A(5))_2&=-44045680773437284570849713497906508911110782976,
 \end{aligned}                                               \tag{2}
\]
and $(\LL^3 A(m))_2<0$ for every integer $m\ge6$.
\end{theorem}

\subsection*{The finite certificate at $m=4$}
Substitution in (1) gives
\[
 (A_1(4),\ldots,A_{10}(4))=
 (1,2,7,32,177,1122,7898,60398,494078,4274228).
                                                               \tag{3}
\]
These are the known terms of OEIS A005362, with index shifted
by one; their numerical discovery is not claimed here.
Applying the displayed definition of $\LL$ gives
\[
\begin{array}{c|r|r}
n&(\LL A(4))_n&(\LL^2 A(4))_n\\\hline
1&3&420\\
2&15&22400\\
3&215&8967705\\
4&4575&5312555556\\
5&139062&6332625182576\\
6&5388152&10985704681554720\\
7&254309640&\\
8&14041752660&
\end{array}                                                  \tag{4}
\]
The next row is
\[
\begin{array}{c|r}
n&(\LL^3 A(4))_n\\\hline
1&3264676100\\
2&38581511487375\\
3&28565867977326238944\\
4&18260024739573021120028544
\end{array}                                                  \tag{5}
\]
Consequently $(\LL^4 A(4))_2$ equals
\[
 38581511487375\cdot18260024739573021120028544
 -28565867977326238944^2,
\]
which is the negative integer in (2). This already disproves
the stated universal conjecture. The entry uses only $A_2$ through
$A_{10}$, so there is no choice of a boundary value involved.

For $m=5$, the ten initial values from (1) are
\[
 (1,2,8,44,310,2606,25202,272582,3233738,41454272).
                                                               \tag{6}
\]
Four applications of the same subtraction give the second value
in (2). The complete intermediate rows for both parameters are
included in the exact verification output.

\subsection*{A polynomial certificate for every $m\ge6$}
Writing $\ell=n-1-k$, factorial cancellation in (1) gives
\[
 A_n(m)=\sum_{k=0}^{n-1}M(k,n-1-k;m),\qquad
 M(k,\ell;m)=\prod_{i=1}^k\prod_{j=1}^{\ell}
                \frac{m+i+j-1}{i+j-1}.                       \tag{7}
\]
For example, both expressions for each summand equal
$\prod_{i=0}^{m-1}(k+\ell+i)!i!/((k+i)!(\ell+i)!)$.
Empty products are $1$. Formula (7) expresses every $A_n(m)$
as a polynomial in $m$ with rational coefficients.

Set
\[
 Q(m)=m^4+9m^3+26m^2+36m+72
\]
and
\[
\begin{aligned}
P(m)={}&m^{17}+42m^{16}+723m^{15}+6462m^{14}+27742m^{13}\\
 &-16148m^{12}-934842m^{11}-5962420m^{10}-21758691m^9\\
 &-53818158m^8-97778329m^7-138523978m^6\\
 &-159475276m^5-154920392m^4-129902448m^3\\
 &-88784352m^2-37081152m-2892672.
\end{aligned}                                                \tag{8}
\]
Expanding (7) for $2\le n\le8$ and applying $\LL$ three times
yields the polynomial identity
\[
 (\LL^3 A(m))_2=
 -\frac{m^4(m-1)^4(m+1)^5(m+2)^2Q(m)P(m)}{1547934105600}.
                                                               \tag{9}
\]
This is a finite algebraic identity, directly reproducible from
(7). The verifier checks it by exact polynomial arithmetic and
independently by 37 evaluations of the source binomial sums;
the degree is at most 36.

It remains to prove the sign uniformly. The function
\[
 \frac{P(m)}{m^{13}}
 =m^4+42m^3+723m^2+6462m+27742
   -\sum_{j=1}^{13}c_jm^{-j}
\]
has $c_j>0$ by (8), so its derivative is strictly positive
for $m>0$. Direct evaluation gives
\[
 P(6)=264264378826752>0.
\]
It follows that $P(m)>0$ for every real $m\ge6$.
All other factors in (9) are positive there, proving the
negative sign and completing the theorem.

\section*{Verification, scope and attribution}
The verifier uses exact integer and rational arithmetic throughout.
It compares the two defining products, reproduces both finite
certificates, checks the degree-36 identity in two ways, and
supplies an additional sign certificate: all 18 coefficients
of $P(z+6)$ are strictly positive.
The infinite family follows from (9) and the proved sign, not
from numerical extrapolation. Review is by the authoring agent only.

The source already establishes ordinary log-convexity at $m=4$;
the negative fourth iterate does not contradict that weaker result.
The Catalan case $m=2$ is known to be infinitely log-convex.
Conjecture 4.2, the Baxter case $m=3$, remains separate.
The row and column log-concavity results in \cite{MS,LM} concern
different sequences. The complete current preprint \cite{FZZ}
was read. Its identified journal final, DOI
\href{https://doi.org/10.1016/j.jmaa.2026.130615}{10.1016/j.jmaa.2026.130615},
remains unread, so this package receives zero new campaign credit
pending the final-source audit.

\section*{Tasks 2--5. Follow-up directions}
\begin{enumerate}[start=2,leftmargin=*]
\item Decide infinite log-convexity for the remaining Baxter
parameter $m=3$.
\item Determine the largest $r$ for which the entire Hoggatt
sequence is $r$-log-convex, as a function of $m\ge4$.
\item Prove eventual $r$-log-convexity for every fixed $m,r$,
with explicit bounds on the starting index.
\item Locate all negative entries of the third and fourth
iterates, and determine how their positions depend on $m$.
\end{enumerate}

\begin{thebibliography}{9}
\raggedright
\bibitem{FZZ} H. Fang, C. X. T. Zhang and J. J. Y. Zhao,
\emph{Analytic properties arising from the Baxter numbers},
\href{https://arxiv.org/abs/2505.05873}{arXiv:2505.05873v1},
current 14-page preprint; Conjecture 4.3 on p.~11.
\bibitem{OEIS} OEIS Foundation,
\href{https://oeis.org/A005362}{A005362}, Hoggatt sequence with parameter $4$.
Entry by N. J. A. Sloane and S. Plouffe.
\bibitem{MS} J. Mao and W. Shi,
\emph{The analytic properties of Hoggatt triangles},
\href{https://arxiv.org/abs/2607.01582}{arXiv:2607.01582v1}, 2026.
\bibitem{LM} Y. Liu and J. Mao,
\emph{Pfaffian--Toeplitz identities, Schur positivity, and the
$q$-log-convexity of Baxter polynomials},
\href{https://arxiv.org/abs/2608.22258}{arXiv:2608.22258v1}, 2026.
\bibitem{DB} MathDB, \href{https://mathdb.com/p/368345}{Problem 368345},
accessed September 18, 2026.
\end{thebibliography}
\end{document}
